% !TEX root = ../main.tex
% ESEMPIO 2: intervalli di overlay non contigui e titolo che cambia a ogni passo.
%   - Un elemento può essere evidenziato in più momenti separati: <1-1, 3-3> (liste separate da virgola).
%     Vale per \alert<...>, \item<...> e per gli stili TikZ: on slide={<1-1, 3-3> highlight light}.
%   - Gli intervalli si possono tenere in una macro e riusare: \alert<\rangeFirstAndThird>{...}.
%   - \only<n>{\frametitle{...}}: il titolo cambia da un passo all'altro.

\begin{frame}
    \only<1>{\frametitle{Example: title of step 1}}
    \only<2->{\frametitle{Example: the title changes with the step}}

    \newcommand{\rangeFirstAndThird}{1-1, 3-3}
    \newcommand{\rangeSecond}{2-2}

    \begin{itemize}\setlength\itemsep{1em}
        \item \alert<\rangeFirstAndThird>{highlighted at steps 1 and 3, range kept in a macro}
        \item \alert<\rangeSecond>{highlighted only at step 2}
        \item \alert<1-1, 3-3>{the same two ranges, written directly}
    \end{itemize}

    \vspace{1em}

    \begin{center}
        \begin{tikzpicture}
            \node[rboxes, on slide={<1-1, 3-3> highlight light}] {a TikZ node highlighted at steps 1 and 3};
        \end{tikzpicture}
    \end{center}

\end{frame}
